\documentclass[10pt,a4paper]{amsart}
\usepackage{amsmath,amssymb,amsthm,graphicx}
% Retained minimal preamble for the Stacks Project chapter flat.tex (More on Flatness).
% Verbatim excerpt of the chapter preamble: Xy-pic setup (native preamble lines 14-18) and the native
% theorem declarations with the mathematical macros (native preamble lines 154-217). Copied mechanically;
% nothing added, removed or redefined. The wrapper class and the amsmath/amssymb/amsthm/graphicx base
% packages are supplied by the assembler, so the class/externaldocument part of the native preamble is omitted.
% For commutative diagrams we use Xy-pic
\usepackage[all]{xy}

% We use 2cell for 2-commutative diagrams.
\xyoption{2cell}
% Theorem environments.
%
\theoremstyle{plain}
\newtheorem{theorem}[subsection]{Theorem}
\newtheorem{proposition}[subsection]{Proposition}
\newtheorem{lemma}[subsection]{Lemma}

\theoremstyle{definition}
\newtheorem{definition}[subsection]{Definition}
\newtheorem{example}[subsection]{Example}
\newtheorem{exercise}[subsection]{Exercise}
\newtheorem{situation}[subsection]{Situation}

\theoremstyle{remark}
\newtheorem{remark}[subsection]{Remark}
\newtheorem{remarks}[subsection]{Remarks}

\numberwithin{equation}{subsection}

% Macros
%
\def\lim{\mathop{\mathrm{lim}}\nolimits}
\def\colim{\mathop{\mathrm{colim}}\nolimits}
\def\Spec{\mathop{\mathrm{Spec}}}
\def\Hom{\mathop{\mathrm{Hom}}\nolimits}
\def\Ext{\mathop{\mathrm{Ext}}\nolimits}
\def\SheafHom{\mathop{\mathcal{H}\!\mathit{om}}\nolimits}
\def\SheafExt{\mathop{\mathcal{E}\!\mathit{xt}}\nolimits}
\def\Sch{\mathit{Sch}}
\def\Mor{\mathop{\mathrm{Mor}}\nolimits}
\def\Ob{\mathop{\mathrm{Ob}}\nolimits}
\def\Sh{\mathop{\mathit{Sh}}\nolimits}
\def\NL{\mathop{N\!L}\nolimits}
\def\CH{\mathop{\mathrm{CH}}\nolimits}
\def\proetale{{pro\text{-}\acute{e}tale}}
\def\etale{{\acute{e}tale}}
\def\QCoh{\mathit{QCoh}}
\def\Ker{\mathop{\mathrm{Ker}}}
\def\Im{\mathop{\mathrm{Im}}}
\def\Coker{\mathop{\mathrm{Coker}}}
\def\Coim{\mathop{\mathrm{Coim}}}

% Boxtimes
%
\DeclareMathSymbol{\boxtimes}{\mathbin}{AMSa}{"02}

%
% Macros for moduli stacks/spaces
%
\def\QCohstack{\mathcal{QC}\!\mathit{oh}}
\def\Cohstack{\mathcal{C}\!\mathit{oh}}
\def\Spacesstack{\mathcal{S}\!\mathit{paces}}
\def\Quotfunctor{\mathrm{Quot}}
\def\Hilbfunctor{\mathrm{Hilb}}
\def\Curvesstack{\mathcal{C}\!\mathit{urves}}
\def\Polarizedstack{\mathcal{P}\!\mathit{olarized}}
\def\Complexesstack{\mathcal{C}\!\mathit{omplexes}}
% \Pic is the operator that assigns to X its picard group, usage \Pic(X)
% \Picardstack_{X/B} denotes the Picard stack of X over B
% \Picardfunctor_{X/B} denotes the Picard functor of X over B
\def\Pic{\mathop{\mathrm{Pic}}\nolimits}
\def\Picardstack{\mathcal{P}\!\mathit{ic}}
\def\Picardfunctor{\mathrm{Pic}}
\def\Deformationcategory{\mathcal{D}\!\mathit{ef}}

% --- appended 2026-10-12 (Chief): the project own definitions of the reference/slogan/history
% display environments, copied mechanically from the shared Stacks preamble lines 8-12
% (\usepackage{verbatim} plus the three \newenvironment lines, which the chapter preamble keeps
% verbatim from the project preamble). Nothing added, removed or redefined: these environments
% carry no mathematical content and the project itself typesets them as comments. Without these
% three lines a counted statement whose own text uses them fails with Environment slogan undefined.
\usepackage{verbatim}
\newenvironment{reference}{\comment}{\endcomment}
\newenvironment{slogan}{\comment}{\endcomment}
\newenvironment{history}{\comment}{\endcomment}
\usepackage{cleveref}
\crefname{originallemma}{original lemma}{original lemmas}
\Crefname{originallemma}{Original lemma}{Original lemmas}
\makeatletter
\newlabel{perfect-lemma-descend-relatively-perfect}{{0DI8}{source}{}{sourcepointer-1}{}}
\newlabel{perfect-lemma-descend-relatively-perfect@cref}{{[originallemma][0][]0DI8}{[1][0][]source}{}{}{}}
\makeatother
\crefname{originallemma}{original lemma}{original lemmas}
\Crefname{originallemma}{Original lemma}{Original lemmas}
\makeatletter
\newlabel{lemma-blowup-complex}{{0ESQ}{source}{}{sourcepointer-2}{}}
\newlabel{lemma-blowup-complex@cref}{{[originallemma][0][]0ESQ}{[1][0][]source}{}{}{}}
\makeatother
\title{On an integral scheme a perfect complex becomes locally free with constant ranks on an open, and of tor dimension at most one after an admissible blowup}
\author{The Stacks Project Authors}
\date{stacks-project:flat.tex:local-snapshot-20261011; GNU FDL 1.2 or later}
\begin{document}
\maketitle
\noindent\textit{Editorial note.} All delivered mathematical text is the author's own, copied from the pinned Stacks Project chapter file flat.tex; no statement or proof step was rewritten, reordered or removed.\par\medskip
\noindent\textit{Editorial note.} The preamble is the chapter's own retained minimal preamble (sources/stacks-flat/retained-preamble.tex), a verbatim excerpt of the native Xy-pic setup and of the native theorem declarations and mathematical macros, followed by the project own definitions of the reference, slogan and history display environments (verbatim package plus three newenvironment lines) that the retained file carries at its end and that the project itself typesets as comments. No macro was invented or redefined, and the wrapper class and the amsmath/amssymb/amsthm/graphicx base packages are supplied by the assembler.\par\medskip
\noindent\textit{Editorial note.} The author opens the delivered proof with a rhetorical appeal to the reader; the sentence is kept verbatim in the text (source lines 13116-13117): "We strongly urge the reader to find their own proof of the existence of U." The argument itself is then supplied in full inside the counted span.\par\medskip
\noindent\textit{Editorial note.} Canonical location of the counted proposition: Stacks tag 0ESR (sources/stacks-flat/pinned-original-source.tex line 13106).\par\medskip
\noindent\textit{Allowed standard prerequisites (editorial).}\begin{itemize}
\item Descent of relatively perfect objects: isomorphic after a point implies isomorphic on a neighbourhood. Derived Categories of Schemes, Lemma labelled perfect-lemma-descend-relatively-perfect (Stacks tag 0DI8), invoked at proof lines 13124-13126. Supplies the open neighbourhood U on which the model complex E' and E agree; the construction of E' from the dimensions at eta is inside the counted span.
\item The admissible blowup making the cohomology sheaves of a perfect complex of tor dimension at most one. Same chapter, Lemma labelled lemma-blowup-complex (source line 13036), invoked at proof lines 13127-13128. Supplies the blowup once a nonempty open with locally free constant rank cohomology is available; the reduction to it is in the counted span.
\end{itemize}\medskip
\ifdefined\hypertarget\hypertarget{sourcepointer-1}{}\fi
\noindent\textit{Source reference.} Original lemma 0DI8: sources/stacks-chow/tags.txt, source lines 16933--16933. Cross-chapter standard prerequisite located in the project tag list staged in the local snapshot; invoked at proof line 13125 of the counted proof. It is retained only as a navigation pointer; every load-bearing step is inside the delivered span.\par\medskip
\ifdefined\hypertarget\hypertarget{sourcepointer-2}{}\fi
\noindent\textit{Source reference.} Original lemma 0ESQ: sources/stacks-flat/pinned-original-source.tex, source lines 13036--13036. Same-chapter pointer into the author environment lemma at source line 13036; invoked at proof line 13128 of the counted proof. It is retained only as a navigation pointer; every load-bearing step is inside the delivered span.\par\medskip
\begin{lemma}
\label{lemma-blowup-complex-integral}
Let $X$ be an integral scheme. Let $E \in D(\mathcal{O}_X)$ be perfect.
Then there exists a nonempty open $U \subset X$
such that $H^i(E|_U)$ is finite locally free of constant rank $r_i$
for all $i \in \mathbf{Z}$ and there exists a $U$-admissible blowup
$b : X' \to X$ such that $H^i(Lb^*E)$ is a perfect
$\mathcal{O}_{X'}$-module of tor dimension $\leq 1$ for all $i \in \mathbf{Z}$.
\end{lemma}

\begin{proof}
We strongly urge the reader to find their own proof of the
existence of $U$. Let $\eta \in X$ be the generic point.
The restriction of $E$ to $\eta$ is isomorphic in $D(\kappa(\eta))$ to a
finite complex $V^\bullet$ of finite dimensional vector spaces with zero
differentials. Set $r_i = \dim_{\kappa(\eta)} V^i$.
Then the perfect object $E'$ in $D(\mathcal{O}_X)$ represented by the complex
with terms $\mathcal{O}_X^{\oplus r_i}$ and zero differentials
becomes isomorphic to $E$ after pulling back to $\eta$.
Hence by Derived Categories of Schemes, Lemma
\ref{perfect-lemma-descend-relatively-perfect}
there is an open neighbourhood $U$ of $\eta$ such that $E|_U$ and $E'|_U$
are isomorphic. This proves the first assertion. The second follows from the
first and Lemma \ref{lemma-blowup-complex} as any nonempty open is
scheme theoretically dense in the integral scheme $X$.
\end{proof}

\end{document}
